% Mẫu một đường: thay metadata, công thức, miền, cửa sổ, điểm và nhãn tại đây.
% Mục đích mẫu: tiếp tuyến nằm ngang không đồng nghĩa với cực trị.
\documentclass[tikz,border=3pt]{standalone}
\usepackage{fontspec}
\usepackage{unicode-math}
\usepackage{pgfplots}
\pgfplotsset{compat=1.18}
\usetikzlibrary{backgrounds,arrows.meta}
\input{zo_graph_v02.tex}

\begin{document}
\begin{tikzpicture}[zo graph v02 field]
\zoGraphVTwoLayers
\begin{axis}[
  zo graph v02 axis,
  xmin=-1.4,xmax=3.4,ymin=-3,ymax=8,
  restrict y to domain=-3:8,
  xtick={-1,1,2,3},ytick={-2,2,4,6},
  xlabel={$x$},ylabel={$y$},
  every axis x label/.append style={at={(axis description cs:0.985,0.31)},anchor=south east},
  every axis y label/.append style={at={(axis description cs:0.31,0.985)},anchor=north west}
]
  % Bảo đảm cả STIX Two Text và STIX Two Math được nhúng.
  \node[zo graph v02 font sentinel] at (axis cs:0,0) {x};

  % Yếu tố phụ: dùng style theo vai trò toán học, không dùng màu đường cong khác.
  \addplot[zo graph v02 tangent,domain=-0.35:2.75,samples=2,forget plot] {2};

  % Đường cong chính.
  \addplot[zo graph v02 curve solid,domain=-1.4:3.4,samples=401,forget plot] {(x-1)^3+2};

  % Marker được vẽ trên đường cong.
  \addplot[zo graph v02 point emphasized,forget plot] coordinates {(1,2)};

  \begin{pgfonlayer}{zoGraphVTwoTextLayer}
    % Nhãn điểm cách marker khoảng 2--3 mm; nền nhỏ không che vùng tiếp xúc.
    \node[zo graph v02 protected label,anchor=north west]
      at (axis cs:1.14,1.72) {$A(1;2)$};
    \node[zo graph v02 protected label,anchor=east]
      at (axis cs:2.78,7.62) {$y=(x-1)^3+2$};
    \node[zo graph v02 protected label,anchor=south east]
      at (axis cs:2.70,2.18) {Tiếp tuyến $y=2$};
    \node[zo graph v02 tick label,anchor=north east,xshift=-2pt,yshift=-2pt]
      at (axis cs:0,0) {$O$};
  \end{pgfonlayer}
\end{axis}
\end{tikzpicture}
\end{document}
